% =============================================================================
% Swaps de taux vanilles (IRS) EUR contre Euribor : theorie, pratique, implementation
% Compiler avec :  pdflatex pricing_irs.tex  (deux passes)
% Les chiffres, tableaux et donnees des figures sont generes par
% generate_irs_data.py a partir du code du depot (voir l'annexe).
% =============================================================================
\documentclass[11pt,a4paper]{article}

\usepackage[utf8]{inputenc}
\usepackage[T1]{fontenc}
\usepackage{lmodern}
\usepackage[french]{babel}
\usepackage{textcomp}
\usepackage{amsmath,amssymb,amsthm,mathtools}
\usepackage{microtype}
\usepackage[a4paper,margin=2.5cm]{geometry}
\usepackage{booktabs,array,tabularx,longtable}
\usepackage{graphicx}
\usepackage{float}
\usepackage[dvipsnames]{xcolor}
\usepackage{enumitem}
\usepackage{listings}
\usepackage{tikz}
\usetikzlibrary{arrows.meta,positioning,shapes.geometric,fit,calc,babel}
\usepackage{pgfplots}
\usepgfplotslibrary{groupplots}
\pgfplotsset{compat=newest}
\pgfkeys{/pgf/number format/.cd, use comma}
\usepackage{fancyhdr}
\usepackage[font=small,labelfont=bf]{caption}
\usepackage[colorlinks=true,linkcolor=NavyBlue,urlcolor=NavyBlue,citecolor=NavyBlue]{hyperref}

% --- Mise en page ------------------------------------------------------------
\setlength{\parindent}{0pt}
\setlength{\parskip}{0.55em}
\setlength{\emergencystretch}{3em}
\setlist{itemsep=0.15em,topsep=0.3em}
\pagestyle{fancy}
\fancyhf{}
\fancyhead[L]{\small IRS EUR}
\fancyhead[R]{\small\leftmark}
\fancyfoot[C]{\small\thepage}
\renewcommand{\headrulewidth}{0.3pt}

% --- Macros ------------------------------------------------------------------
\newcommand{\estr}{\texorpdfstring{\texteuro STR}{EUR STR}}
\newcommand{\code}[1]{\texttt{\detokenize{#1}}}
\newcommand{\pct}[1]{#1\,\%}
\newcommand{\eur}{\ifmmode\,\text{\texteuro}\else\,\texteuro\fi}
\newcommand{\bp}{\,\mathrm{bp}}
\DeclareMathOperator{\PV}{PV}

\theoremstyle{plain}
\newtheorem{proposition}{Proposition}
\theoremstyle{definition}
\newtheorem{definition}{Définition}
\theoremstyle{remark}
\newtheorem*{remarque}{Remarque}

% --- Listings ----------------------------------------------------------------
\definecolor{codebg}{gray}{0.96}
\definecolor{codekw}{rgb}{0.10,0.25,0.55}
\definecolor{codecm}{rgb}{0.30,0.45,0.30}
\lstset{
  language=Python,
  basicstyle=\fontsize{7.6}{9.2}\selectfont\ttfamily,
  keywordstyle=\color{codekw}\bfseries,
  commentstyle=\color{codecm}\itshape,
  stringstyle=\color{BrickRed},
  backgroundcolor=\color{codebg},
  frame=single,
  rulecolor=\color{black!25},
  breaklines=true,
  columns=fullflexible,
  keepspaces=true,
  showstringspaces=false,
  aboveskip=0.6em,
  belowskip=0.4em,
  captionpos=b,
}
\renewcommand{\lstlistingname}{Listing}
\lstdefinelanguage{json}{morestring=[b]", sensitive=false}

\pgfplotsset{
  every axis/.append style={
    grid=major, grid style={black!12}, tick label style={font=\footnotesize},
    label style={font=\footnotesize}, legend style={font=\footnotesize, draw=black!30},
  },
}
\tikzset{
  proc/.style={draw, rounded corners=3pt, minimum width=3.1cm, minimum height=0.95cm,
    align=center, fill=blue!6, font=\small},
  data/.style={draw, rounded corners=3pt, minimum height=0.8cm, align=center,
    fill=orange!12, font=\small},
  flow/.style={-{Stealth[length=2.4mm]}, thick},
}

% --- Chiffres generes par generate_irs_data.py -------------------------------
% Generated by docs/latex/generate_irs_data.py: do not edit by hand.
\newcommand{\IrsPillarCount}{16}
\newcommand{\IrsMaxIterations}{13}
\newcommand{\IrsMaxResidual}{$1{,}7\times10^{-16}$}
\newcommand{\IrsNotional}{10\,000\,000}
\newcommand{\IrsFixedRate}{2{,}50}
\newcommand{\IrsStart}{2026-10-21}
\newcommand{\IrsEnd}{2031-10-21}
\newcommand{\IrsPV}{$-$31\,884{,}20}
\newcommand{\IrsFixedLegPV}{$-$1\,172\,213{,}17}
\newcommand{\IrsFloatingLegPV}{1\,140\,328{,}97}
\newcommand{\IrsParRate}{2{,}4320}
\newcommand{\IrsAnnuity}{46\,888\,526{,}75}
\newcommand{\IrsPVBP}{4\,688{,}85}
\newcommand{\IrsSpread}{$-$6{,}8}
\newcommand{\IrsFirstFixingDate}{2026-10-19}
\newcommand{\IrsFirstFixing}{2{,}014}
\newcommand{\IrsFloatingCount}{10}
\newcommand{\IrsFixedCount}{5}
\newcommand{\IrsSingleCurvePV}{$-$93\,054{,}94}
\newcommand{\IrsSingleCurveGap}{$-$61\,170{,}74}
\newcommand{\IrsParallelDV}{4\,197{,}48}
\newcommand{\IrsConvexity}{$-$2{,}714}
\newcommand{\IrsDVOis}{9{,}39}
\newcommand{\IrsDVSix}{4\,188{,}10}
\newcommand{\IrsDVThree}{0{,}00}
\newcommand{\IrsKeyRateSum}{4\,197{,}48}
\newcommand{\IrsKeyRateGap}{$3{,}2\times10^{-4}$}
\newcommand{\IrsRiskBumps}{122}
\newcommand{\IrsTaylorShift}{10}
\newcommand{\IrsTaylorActual}{41\,839{,}42}
\newcommand{\IrsTaylorEstimate}{41\,839{,}10}


\title{\textbf{Les swaps de taux vanilles (IRS) en euro}\\[0.4em]
  \large Théorie, pratique de marché et implémentation dans le pricer de dérivés de taux}
\author{Projet \texttt{rate-derivatives-pricer} -- phase 2}
\date{\today}

\begin{document}
\maketitle
\thispagestyle{empty}

\begin{abstract}
\noindent
Ce document présente le swap de taux d'intérêt vanille (IRS) contre Euribor : sa définition, sa
valorisation dans le cadre \emph{multi-courbes} utilisé depuis la crise de 2008 (actualisation sur
la courbe \estr{}, projection sur la courbe de l'indice), les conventions de marché de l'euro, la
construction des courbes par \emph{bootstrap} et les mesures de risque par choc des cotations
(DV01 parallèle, par instrument de calibration, par courbe, convexité). Pour chaque notion, il montre
comment elle est implémentée dans le pricer : fichier, classe, extrait de code et test. Un exemple
chiffré complet (swap payeur à 5 ans) et la validation contre QuantLib terminent le document. Chaque
nombre issu du pricer est produit par le code du dépôt.

\medskip\noindent
\textbf{Avertissement.} Les données de marché (snapshot \code{EUR_2026-10-19}) sont
\emph{synthétiques} : aucun niveau de taux ou de base de ce document ne doit être lu comme une
cotation réelle.
\end{abstract}

\newpage
{\setlength{\parskip}{0pt}\tableofcontents}
\newpage

% =============================================================================
\part{Théorie et pratique de marché}
% =============================================================================

\section{Le swap de taux vanille}\label{sec:swap}

\subsection{Définition}
Un swap de taux vanille échange, entre deux dates \(T_0 < T_n\), les intérêts d'un taux \emph{fixe}
\(K\) contre ceux d'un taux \emph{variable} indexé sur un taux interbancaire à terme, ici
l'Euribor 6 mois (ou 3 mois), pour un notionnel \(N\) qui n'est jamais échangé.
\begin{itemize}
  \item La \textbf{jambe fixe} verse à chaque date \(T_i^{\mathrm{fix}}\) le montant
    \(N\,K\,\tau_i^{\mathrm{fix}}\), où \(\tau_i^{\mathrm{fix}}\) est la fraction d'année de la période.
  \item La \textbf{jambe flottante} verse à chaque date \(T_j\) le montant \(N\,L_j\,\tau_j\), où
    \(L_j\) est l'Euribor observé (\emph{fixé}) au début de la période \([S_j,T_j]\), avec
    \(S_j = T_{j-1}\) : le taux est connu en début de période et payé en fin de période
    (\emph{in advance}).
\end{itemize}
Le \textbf{payeur} (\code{pay_fixed}) paie le fixe et reçoit le flottant ; le \textbf{receveur}
(\code{receive_fixed}) fait l'inverse. Le nom de la direction désigne toujours la jambe fixe.

\begin{figure}[H]
\centering
\begin{tikzpicture}[x=1cm,y=1cm]
  \draw[-{Stealth}, thick] (0,0) -- (12.2,0) node[right] {\(t\)};
  \foreach \x/\l in {0/{T_0},3/{T_1^{\mathrm{fix}}},6/{T_2^{\mathrm{fix}}},9/{T_3^{\mathrm{fix}}},12/{T_n}}
    \draw (\x,0.1) -- (\x,-0.1) node[below] {\(\l\)};
  \foreach \x in {3,6,9,12}
    \draw[-{Stealth}, thick, blue!70!black] (\x,0) -- (\x,1.1);
  \node[blue!70!black, font=\small] at (7.5,1.45) {jambe fixe : \(N K \tau^{\mathrm{fix}}\), annuelle, 30/360};
  \foreach \x in {1.5,3,4.5,6,7.5,9,10.5,12}
    \draw[-{Stealth}, thick, orange!85!black] (\x,0) -- (\x,-1.3);
  \node[orange!85!black, font=\small] at (7.5,-1.85) {jambe flottante : \(N L_j \tau_j\), semestrielle, Act/360};
\end{tikzpicture}
\caption{Flux d'un swap payeur contre Euribor 6M (schéma). Le payeur reçoit la jambe flottante (en
bas) et paie la jambe fixe (en haut) ; les flux de signes opposés se compensent aux dates communes.}
\end{figure}

\subsection{Pourquoi cet instrument}
L'IRS est l'instrument de base du marché des taux : il transforme une dette à taux variable en dette à
taux fixe (ou l'inverse), couvre un portefeuille obligataire, et ses taux de par par maturité
(\emph{cotations de swaps}) sont l'information à partir de laquelle on construit les courbes de taux.

\subsection{Rappels : facteurs d'actualisation et taux forward}
Pour une courbe \(c\), \(P^c(t)\) est le facteur d'actualisation de la date de valorisation à \(t\)
(\(P^c(0)=1\)). Le taux forward simple entre \(s\) et \(t\), avec la fraction d'année \(\tau(s,t)\) de
l'indice, vaut
\begin{equation}
  F^c(s,t) \;=\; \frac{1}{\tau(s,t)}\left(\frac{P^c(s)}{P^c(t)} - 1\right).
  \label{eq:forward}
\end{equation}
Le document \texttt{pricing\_ois.pdf} détaille ces notions, les décomptes de jours et le calendrier
TARGET ; elles sont réutilisées telles quelles.

\section{Valorisation : du mono-courbe au multi-courbes}\label{sec:pricing}

\subsection{L'ancien cadre à courbe unique}
Avant 2008, on supposait que l'Euribor se projetait sur la même courbe que celle qui actualise. En
notant \(P\) cette courbe, le forward vaut \(F_j = F(S_j,T_j)\) et le flux flottant actualisé
\(N\,\tau_j F_j P(T_j) = N\,(P(S_j)-P(T_j))\) : la somme \emph{télescope} en
\begin{equation}
  \PV_{\mathrm{flottant}} \;=\; N\,\bigl(P(T_0) - P(T_n)\bigr),
  \label{eq:telescope}
\end{equation}
si bien que le taux de par du swap ne dépend plus que de la courbe unique.

\subsection{Pourquoi le multi-courbes}
La crise de 2008 a fait apparaître des écarts durables entre l'Euribor à différentes échéances et le
taux sans risque : l'Euribor 6M embarque une prime de crédit et de liquidité interbancaire que l'\estr{}
n'a pas, et l'Euribor 6M est supérieur à l'Euribor 3M. Deux conséquences :
\begin{enumerate}
  \item \textbf{Actualisation.} Les swaps compensés ou collatéralisés sont rémunérés au taux de leur
    collatéral, \estr{} en euro : tous les flux s'actualisent sur la \emph{courbe OIS} \(P^{d}\).
  \item \textbf{Projection.} Chaque indice a sa \emph{propre} courbe de projection \(P^{e}\) (une pour
    l'Euribor 6M, une pour l'Euribor 3M), calibrée sur les swaps cotés contre cet indice.
\end{enumerate}
La relation \eqref{eq:telescope} n'est alors plus vraie : projeter et actualiser sur des courbes
différentes rend le taux de par dépendant des deux courbes.

\subsection{Valeur actuelle}
Pour un payeur, avec \(F^e_j = F^e(S_j,T_j)\) calculé par \eqref{eq:forward} sur la courbe de
projection \(e\) de l'indice :
\begin{align}
  \PV_{\mathrm{flottant}} &= \sum_{j} N\,\tau_j\,F^e_j\,P^d(T_j), \label{eq:floatpv}\\
  \PV_{\mathrm{fixe}}     &= \sum_{i} N\,K\,\tau^{\mathrm{fix}}_i\,P^d(T^{\mathrm{fix}}_i) \;=\; K\,A,
  \label{eq:fixedpv}\\
  \PV_{\mathrm{payeur}}   &= \PV_{\mathrm{flottant}} - \PV_{\mathrm{fixe}}, \qquad
  \PV_{\mathrm{receveur}} = -\PV_{\mathrm{payeur}}. \label{eq:pv}
\end{align}
Les sommes ne portent que sur les coupons \emph{non encore payés} (date de paiement strictement
postérieure à la date de valorisation). Le facteur \(A\) est l'\textbf{annuité}.

\subsection{Taux de par, annuité et PVBP}
\begin{definition}
L'\emph{annuité} est \(A=\sum_i N\,\tau_i^{\mathrm{fix}}\,P^d(T_i^{\mathrm{fix}})\) ; le \emph{taux de par}
est le taux fixe qui annule la valeur actuelle,
\begin{equation}
  S \;=\; \frac{\PV_{\mathrm{flottant}}}{A},
  \qquad\text{et}\qquad
  \PV_{\mathrm{payeur}} \;=\; A\,(S-K).
  \label{eq:par}
\end{equation}
Le \emph{PVBP} est la variation de valeur d'un point de base du taux fixe, \(A\times10^{-4}\).
\end{definition}
Le taux de par \(S\) est la cotation de marché d'un swap de la même maturité (\S\ref{sec:bootstrap}).
Un swap dont le taux fixe est le par rate a une valeur actuelle nulle, c'est aussi la propriété testée
par le pricer.

\subsection{Le coupon indexé sur Euribor}\label{sec:iborcoupon}
Le taux du coupon qui commence en \(S_j\) est \emph{fixé} à la \textbf{date de fixing}
\(f_j\), deux jours ouvrés TARGET avant \(S_j\). Selon la position de \(f_j\) par rapport à la date de
valorisation \(t_0\) :
\begin{center}
\begin{tabular}{@{}lll@{}}
\toprule
Cas & Taux \(L_j\) du coupon & Si le fixing manque \\
\midrule
\(f_j < t_0\) (fixing passé) & fixing publié & erreur (\code{MarketDataError}) \\
\(f_j = t_0\) (fixing du jour) & fixing publié s'il est fourni & projeté \\
\(f_j > t_0\) & forward \(F^e(S_j,T_j)\) de la courbe de projection & --- \\
\bottomrule
\end{tabular}
\end{center}
Le forward projeté est lu \emph{sur la période d'accrual elle-même} : c'est la convention dite
\emph{par coupon} de QuantLib. Un fixing connu ne dépend plus du marché : c'est ce qui rend le premier
coupon d'un swap neuf insensible aux chocs de taux (\S\ref{sec:risk}).

\section{Conventions de marché de l'euro}\label{sec:conv}

\begin{center}
\begin{tabularx}{\linewidth}{@{}lX@{}}
\toprule
Élément & Convention retenue (ADR-027, fichier \code{EUR.json}) \\
\midrule
Calendrier & TARGET \\
Date spot & date de valorisation (ajustée \emph{following}) + 2 jours ouvrés \\
Jambe fixe & annuelle, \textbf{30/360 Bond Basis} (ISDA 4.16(f)) \\
Jambe flottante & fréquence de l'indice (6M ou 3M), Act/360 \\
Indice & Euribor 6M (défaut) ou 3M ; fixing 2 jours ouvrés TARGET avant le début, Act/360 \\
Ajustement des dates & Modified Following, y compris l'échéance \\
Génération & \emph{backward} depuis l'échéance, stub éventuel en tête ; règle de fin de mois désactivée \\
Décalage de paiement & aucun (paiement à la date de fin de période) \\
Échéance & départ + tenor, puis ajustée par l'échéancier \\
Dépôt Euribor (calibration) & départ spot, échéance = tenor de l'indice, Act/360, Modified Following, fin de mois \\
\bottomrule
\end{tabularx}
\end{center}
Ces conventions sont celles de \code{EuriborSwapIsdaFixA} de QuantLib. Elles ne se devinent pas : elles
ont été validées avant d'être écrites (règle 7 du projet), puis vérifiées date par date contre QuantLib.

\section{Construction des courbes par bootstrap multi-courbes}\label{sec:bootstrap}

\subsection{Instruments de calibration}
\begin{itemize}
  \item \textbf{Courbe \estr{} (actualisation)} : un dépôt overnight puis des OIS de 1 semaine à
    30 ans (phase~1).
  \item \textbf{Courbe de chaque indice} : un \emph{dépôt Euribor} (spot, durée du tenor de l'indice),
    puis des IRS annuels de 1 à 30 ans contre cet indice. Pas de FRA ni de \emph{basis swap} dans cette
    phase.
\end{itemize}
Chaque instrument a une \emph{cotation implicite} que l'on sait calculer à partir de courbes
candidates, et un \emph{pilier} : la dernière date dont sa valeur dépend (la dernière date de paiement).

\subsection{Algorithme en deux étapes}
\begin{enumerate}
  \item Calibrer la courbe \(P^d\) (\estr{}) sur ses dépôts et OIS, comme en phase 1.
  \item Pour chaque indice \(e\) : trier ses instruments par pilier ; pour chaque pilier \(n\),
    trouver le facteur \(P^e_n\) tel que la cotation implicite égale la cotation de marché, les nœuds
    \(P^e_1,\dots,P^e_{n-1}\) étant déjà connus et \(P^d\) \emph{fixée}.
\end{enumerate}
Pour un IRS de pilier \(T_n\), l'équation est
\begin{equation}
  \frac{1}{A}\sum_{j} N\,\tau_j\,F^e_j(P^e_n)\,P^d(T_j) \;=\; S^{\mathrm{marché}}_n,
  \label{eq:calib-swap}
\end{equation}
et pour le dépôt, \(F^e(\mathrm{spot},T)=q^{\mathrm{marché}}\). Avec une interpolation locale
(\emph{log-linéaire sur les facteurs d'actualisation}), l'instrument \(n\) ne dépend que des nœuds
jusqu'à \(n\) : une passe séquentielle est exacte, et la recherche de racine est unidimensionnelle
(Brent).

\subsection{Calibrer sans les fixings}\label{sec:nofixing}
Un IRS spot a pour premier fixing celui de la date de valorisation. Si cette valeur publiée était
visible pendant la calibration, le premier coupon serait \emph{figé} au lieu d'être projeté et la
courbe dépendrait d'un fixing, ce qu'une cotation de marché ne dit pas. Les instruments de calibration
sont donc valorisés \emph{sans fixings} ; le swap que l'on valorise ensuite, lui, les utilise.

\section{Risque par choc des cotations}\label{sec:risk}

\subsection{Principe : bump-and-reprice sur les instruments de calibration}
On choque de \(h\) points de base le taux coté d'instruments de calibration, on \emph{re-calibre}
toutes les courbes, on revalorise. Le résultat est la sensibilité aux \emph{par rates}, celle que
regardent les desks, par opposition à un choc direct des taux zéro. Soit \(\PV(\mathbf q)\) la valeur
en fonction du vecteur des cotations.

\begin{center}
\begin{tabularx}{\linewidth}{@{}lX@{}}
\toprule
Mesure & Définition (ADR-029) \\
\midrule
DV01 parallèle & \(\bigl[\PV(\mathbf q+h\mathbf 1) - \PV(\mathbf q-h\mathbf 1)\bigr]/2h\), par bp \\
DV01 par instrument (\emph{key rate}) & même formule en ne choquant qu'une cotation \(q_k\) à la fois \\
DV01 par courbe & on choque les cotations d'une courbe : OIS (actualisation, avec re-calibration des
courbes Ibor) ou un indice Ibor (projection) \\
Convexité & \(\bigl[\PV(+h) - 2\PV + \PV(-h)\bigr]/h^2\), en bp\(^2\), choc parallèle \\
\bottomrule
\end{tabularx}
\end{center}
Le signe est celui du PV : le DV01 est la variation de valeur pour \(+1\bp\) ; un payeur de fixe gagne
quand les taux montent (DV01 \(>0\)).

\subsection{Propriétés utiles}
\begin{itemize}
  \item \textbf{Différence centrée.} Elle annule les termes d'ordre deux : l'erreur est d'ordre
    \(h^2\). Une différence avant ou arrière (\code{up}, \code{down}) y ajoute \(\pm\tfrac12 h\,\PV''\).
  \item \textbf{Cohérence.} La somme des DV01 par instrument, comme celle des DV01 par courbe, vaut le
    DV01 parallèle (dérivée d'une somme de chocs), à l'erreur de la différence près.
  \item \textbf{Développement de Taylor.} Pour un choc parallèle \(x\) en points de base,
    \(\PV(x)\approx \PV + \mathrm{DV01}\,x + \tfrac12\,\mathrm{convexité}\,x^2\).
  \item \textbf{Fixings connus.} Un fixing publié ne bouge pas avec le marché : le premier coupon d'un
    swap neuf ne contribue pas au DV01.
\end{itemize}

% =============================================================================
\part{Implémentation dans le pricer}
% =============================================================================

\section{Carte du code}\label{sec:map}

\begin{center}
\begin{tabularx}{\linewidth}{@{}>{\raggedright\arraybackslash}p{4.2cm}X@{}}
\toprule
Notion & Où elle vit (\code{src/rate_derivatives_pricer/}) \\
\midrule
Conventions EUR (config.) & \code{core/conventions/data/EUR.json}, \code{conventions/currency.py} \\
Indice Euribor & \code{core/market/indices.py} : \code{IborIndex} \\
Cotations Ibor & \code{core/market/quotes.py} : \code{IborDepositQuote}, \code{IborSwapQuote} \\
Snapshot, fixings & \code{core/market/snapshot.py}, \code{providers.py} \\
Échéancier & \code{core/schedule/generator.py} (réutilisé de la phase 1) \\
Coupon Euribor & \code{core/instruments/cashflows.py} : \code{IborCoupon} \\
Swap vanille & \code{core/instruments/swaps.py} : \code{VanillaIrs} ; fabrique \code{factories.py} \\
Moteur de valorisation & \code{core/engines/discounting.py} : \code{DiscountingSwapEngine} \\
Instruments de calibration & \code{core/market/calibration.py} \\
Bootstrap multi-courbes & \code{core/market/bootstrap.py} : \code{CurveBuilder} \\
Risque & \code{core/risk/} : \code{bumps.py}, \code{measures.py}, \code{engine.py} \\
Service, API, SDK, UI & \code{services/}, \code{api/}, \code{sdk/}, \code{ui/} \\
\bottomrule
\end{tabularx}
\end{center}
La règle de dépendance reste \(\texttt{ui}\to\texttt{sdk}\to\texttt{api}\to\texttt{services}\to
\texttt{core}\), vérifiée par \code{import-linter} ; le cœur n'importe ni FastAPI ni QuantLib.

\section{Conventions, indices et données de marché}

\subsection{Les conventions sont de la configuration}
Aucune convention n'est codée en dur : la section \code{ibor} de \code{EUR.json} décrit les deux
indices, le dépôt et le swap. Ajouter l'USD ou changer une convention revient à modifier un fichier.
\begin{lstlisting}[language=json,caption={Extrait de \code{EUR.json} : la section \code{ibor}.}]
"ibor": {
  "default_index": "EURIBOR6M",
  "indices": [{"name": "EURIBOR6M", "tenor": "6M", "fixing_days": 2, "calendar": "TARGET",
               "day_count": "ACT/360", "business_day_convention": "MODIFIED_FOLLOWING",
               "end_of_month": true}, ... ],
  "deposit": {"settlement_days": 2, "business_day_convention": "MODIFIED_FOLLOWING",
              "end_of_month": true, "day_count": "ACT/360"},
  "swap": {"fixed_frequency": "ANNUAL", "fixed_day_count": "30/360",
           "business_day_convention": "MODIFIED_FOLLOWING", "end_of_month": false,
           "rule": "BACKWARD", "payment_lag_days": 0}
}
\end{lstlisting}
\code{CurrencyConventions} charge cette section dans des objets figés (\code{IborIndex},
\code{IborSwapConventions}) ; une devise sans section \code{ibor} reste valide et refuse simplement les
IRS (\code{ConventionError}).

\subsection{\code{IborIndex} et date de fixing}
\code{IborIndex} est un objet valeur (dataclass figée) : nom, tenor, jours de fixing, calendrier,
décompte. Sa méthode \code{fixing_date(accrual_start)} recule de \code{fixing_days} jours ouvrés du
calendrier de fixing.

\subsection{Cotations Ibor}
Une cotation Ibor nomme l'indice dont elle calibre la courbe ; il fait partie de son identité, de sorte
que « 5Y » existe une fois par courbe. Le champ \code{index} est obligatoire pour les types
\code{ibor_deposit} et \code{ibor_swap}, interdit pour les autres.
\begin{lstlisting}[caption={\code{IborQuote} : l'indice fait partie de la clé d'identité.}]
@dataclass(frozen=True, slots=True)
class IborQuote(Quote):
    index_name: str = field(kw_only=True)

    @property
    def key(self) -> tuple[str, str]:
        return (self.instrument_type, f"{self.index_name}/{self.tenor_label.upper()}")
\end{lstlisting}
Chaque cotation porte une \emph{unité explicite} (\code{percent}, \code{decimal}, \code{bp}), source
classique de bugs en production ; \code{with_rate()} renvoie une copie dans la même unité, ce qui sert
aux chocs.

\section{Instruments : coupon, swap, échéancier}

\subsection{\code{IborCoupon}}
Le coupon décide lui-même de son taux selon la position de son fixing (tableau de la
\S\ref{sec:iborcoupon}).
\begin{lstlisting}[caption={\code{IborCoupon.rate} : fixing publié ou forward de la courbe de l'indice.}]
def rate(self, env: MarketEnvironment) -> float:
    fixing = env.fixing(self.index.name, self.fixing_date)
    if self.fixing_date < env.valuation_date:
        if fixing is None:
            raise MarketDataError(f"missing {self.index.name} fixing", {...})
        return fixing
    if self.fixing_date == env.valuation_date and fixing is not None:
        return fixing
    return env.projection_curve(self.index.name).forward_rate(
        self.accrual_start, self.accrual_end, self.index.day_count)
\end{lstlisting}
L'environnement de marché (\code{MarketEnvironment}) fournit la courbe de projection \emph{de l'indice},
distincte de la courbe d'actualisation : c'est là que le multi-courbes se joue dans le code.

\subsection{\code{VanillaIrs} et fabrique}
\code{VanillaIrs} (dataclass figée) ne contient que des données : la jambe fixe, la jambe flottante,
l'indice, les dates. Elle ne porte pas son moteur de valorisation. La fabrique
\code{build_vanilla_irs} construit les deux échéanciers avec le même générateur que la phase~1 (fixe :
tenor 1Y ; flottant : tenor de l'indice), puis les coupons.
\begin{lstlisting}[caption={Construction de la jambe flottante (extrait de \code{build_vanilla_irs}).}]
floating_leg = Leg(
    name="floating",
    cashflows=tuple(
        IborCoupon(notional=notional, accrual_start=p.accrual_start, accrual_end=p.accrual_end,
                   payment_date=p.payment_date, day_count=index.day_count, index=index,
                   fixing_date=index.fixing_date(p.accrual_start))
        for p in floating_schedule.periods),
    notional=notional, direction=direction.floating_leg_direction)
\end{lstlisting}
Les dates (début, fin, paiement, fixing) des deux jambes sont identiques à celles de QuantLib
(\S\ref{sec:valid}). La fabrique \code{vanilla_irs_factory} lit une description de trade (JSON) :
devise, indice (par défaut celui de la devise), notionnel, taux fixe, direction, et exactement une de
\code{maturity} ou \code{end_date}, plus une \code{start_date} optionnelle (swap en cours ou forward).

\section{Le moteur de valorisation}

Un seul \code{DiscountingSwapEngine} valorise l'OIS et l'IRS : il somme, jambe par jambe, les flux non
encore payés actualisés sur la courbe d'actualisation, applique le signe de la jambe, et calcule aussi
taux de par, annuité et PVBP (équations \eqref{eq:floatpv}--\eqref{eq:par}). Les moteurs sont sans état
et choisis par un registre ; l'instrument ne connaît pas son moteur.

\subsection{Projection en lot}
Le calcul des taux forward coupon par coupon dominait le temps de calcul d'un bootstrap. Les coupons
exposent donc un point d'extension, \code{Cashflow.amounts(cashflows, env)}, dont la valeur par défaut
boucle sur \code{amount()} ; \code{IborCoupon} le surcharge et lit les facteurs d'actualisation de tous
ses coupons projetés en \emph{un seul appel} à la courbe.
\begin{lstlisting}[caption={Le moteur évalue une jambe par lots, un lot par type de coupon.}]
def _coupon_amounts(coupons, env):
    amounts = [0.0] * len(coupons)
    positions_by_type = {}
    for position, coupon in enumerate(coupons):
        positions_by_type.setdefault(type(coupon), []).append(position)
    for coupon_type, positions in positions_by_type.items():
        batch = coupon_type.amounts([coupons[p] for p in positions], env)
        for position, amount in zip(positions, batch, strict=True):
            amounts[position] = amount
    return amounts
\end{lstlisting}
Le résultat est identique, à l'arrondi près, au calcul coupon par coupon (test unitaire). Le moteur
expose aussi \code{present_value()}, un chemin léger (sans tableau de flux ni journalisation) utilisé
par le risque, qui revalorise des centaines de fois.

\section{Calibration et bootstrap}

\subsection{Instruments de calibration}
\code{CalibrationInstrument} est le pendant du \code{RateHelper} de QuantLib : pour une cotation et un
environnement candidat, il calcule la cotation implicite. Le bootstrap reste générique.
\begin{center}
\begin{tabular}{@{}lll@{}}
\toprule
Classe & Cotation implicite & Pilier \\
\midrule
\code{IborDepositCalibration} & forward simple de la courbe de l'indice, spot \(\to\) tenor & échéance \\
\code{IborSwapCalibration} & taux de par du swap (moteur) & dernier paiement \\
\bottomrule
\end{tabular}
\end{center}

\subsection{\code{CurveBuilder} en deux étapes}
\code{CurveBuilder.build(snapshot)} calibre d'abord la courbe \estr{}, puis une courbe par indice, avec
un environnement de calibration dont la courbe d'actualisation est la courbe \estr{} figée et les
fixings sont \emph{absents} (\S\ref{sec:nofixing}). L'environnement renvoyé contient la courbe
d'actualisation, toutes les courbes de projection, les fixings et un rapport de calibration par
courbe (pilier, cotation, itérations, résidu).

\subsection{Rendre le risque abordable}
Un \emph{key rate} sur 56 cotations demande environ 110 re-bootstraps. Le premier jet prenait 0{,}8~s par
construction, soit 43~s pour un IRS à 5 ans. Après profilage, quatre optimisations \emph{exactes}
(résultats identiques au bit près, testés) ont ramené ce temps à environ 2{,}5~s (6~s pour 30 ans) :
\begin{enumerate}
  \item \textbf{Mémoïsation des piliers résolus.} La clé contient tout ce dont la solution dépend :
    l'instrument, sa cotation, les nœuds déjà résolus de la courbe et, pour un IRS, les nœuds de la
    courbe d'actualisation jusqu'au pilier (plus le nœud qui l'encadre). Un choc ne re-résout que les
    piliers qu'il peut déplacer. Valable pour une interpolation \emph{locale} (attribut
    \code{Interpolator.local}).
  \item \textbf{Mémoïsation des instruments de calibration} : leurs échéanciers ne dépendent pas du
    taux coté.
  \item \textbf{Horizon.} \code{build(snapshot, horizon=d)} s'arrête, pour chaque courbe, au premier
    pilier au moins égal à \(d\) ; le pricer de risque passe la dernière date de flux du trade. Une
    cotation au-delà de l'horizon n'a aucun effet : son bucket est \emph{exactement} nul.
  \item \textbf{Projection en lot} et année-fractions précalculées.
\end{enumerate}
\begin{lstlisting}[caption={Clé de mémoïsation d'un pilier (extrait de \code{_calibrate_pillar}).}]
memo_key = (
    template.name, template.curve_reference_date, instrument.quote.key, pillar, market_quote,
    tuple(pillar_dates), tuple(discount_factors),
    self._discount_signature(discount_curve, pillar),
)
\end{lstlisting}

\section{Le module de risque}

Le risque dépend d'une \emph{fonction} \code{pricer(instrument, snapshot) -> PV}, jamais d'un moteur :
remplacer le produit ou le moteur ne change rien dans \code{risk/}. C'est l'appelant (le service) qui
construit ce pricer : re-bootstrap du snapshot choqué, puis moteur.
\begin{itemize}
  \item \code{BumpTarget} sélectionne les cotations (par courbe, par type, par identité) ;
    \code{BumpSpec} ajoute la taille et le schéma (\code{up}, \code{down}, \code{central}).
  \item \code{Bumper.apply(snapshot, target, shift_bp)} renvoie un \emph{nouveau} snapshot (rien n'est
    modifié) ; une cible qui ne sélectionne aucune cotation est une erreur.
  \item Les mesures (\code{ParallelDV01}, \code{KeyRateDV01}, \code{CurveSpecificDV01},
    \code{Convexity}) sont enregistrées dans un registre ; ajouter une mesure est ajouter un fichier.
    Un test l'enregistre depuis \code{tests/} sans toucher à \code{risk/}.
  \item \code{RiskEngine} mémorise les valorisations par signature de snapshot : le choc parallèle sert
    à la fois le DV01 parallèle et la convexité.
\end{itemize}
\begin{lstlisting}[caption={Première dérivée par bp selon le schéma de différence finie.}]
def dv01(self, spec: BumpSpec) -> float:
    pvs = self.pvs(spec)          # {shift_bp: PV} apres re-bootstrap
    h = spec.size_bp
    match spec.mode:
        case BumpMode.UP:      return (pvs[h] - self.base_pv()) / h
        case BumpMode.DOWN:    return (self.base_pv() - pvs[-h]) / h
        case BumpMode.CENTRAL: return (pvs[h] - pvs[-h]) / (2.0 * h)
\end{lstlisting}

\section{Service, API, SDK et tableau de bord}

\code{PricerService.risk} valide les noms de mesures avant tout calcul, construit le trade, crée un
\code{CurveBuilder} \emph{par requête} (sa mémoire n'est pas partagée) et appelle le moteur de risque.
Les courbes choquées ne passent pas par le cache de courbes. L'API expose \code{POST /v1/risk} ; une
mesure inconnue donne \code{400 UNSUPPORTED_RISK_MEASURE}.
\begin{lstlisting}[language=json,caption={Requête \code{POST /v1/risk} pour un IRS payeur à 5 ans.}]
{ "market":   { "snapshot": "EUR_2026-10-19" },
  "trade":    { "type": "vanilla_irs",
                "params": { "notional": 10000000, "fixed_rate": 0.025,
                            "direction": "pay_fixed", "maturity": "5Y",
                            "index": "EURIBOR6M" } },
  "measures": ["parallel_dv01", "key_rate_dv01", "curve_dv01", "convexity"],
  "bump":     { "size_bp": 1.0, "mode": "central" } }
\end{lstlisting}
Le SDK expose \code{vanilla_irs(...)} et \code{PricerClient.risk(...)} ; le tableau de bord Streamlit
(onglets Marché, Trade, Risque) n'appelle que le SDK. La réponse contient, par mesure, la valeur, l'unité
et les \emph{buckets} (courbe, instrument, valeur).

% =============================================================================
\part{Exemple chiffré et validation}
% =============================================================================

\section{Exemple : un swap payeur à 5 ans contre Euribor 6M}\label{sec:example}

Swap payeur, notionnel \IrsNotional\eur, taux fixe \pct{\IrsFixedRate}, Euribor 6M, du \IrsStart{} au
\IrsEnd{} (départ spot), valorisé au 2026-10-19 sur le snapshot synthétique.

\subsection{Les courbes}
Le bootstrap produit trois courbes : \estr{}, Euribor 3M et Euribor 6M. La
figure~\ref{fig:curves} montre les taux zéro et la \emph{base} Euribor 6M $-$ OIS, mesurée sur le
forward 6M : la courbe 6M est au-dessus de l'OIS aux maturités courtes et moyennes, et celle de 3M entre
les deux. Attention à ne pas sur-interpréter la base à long terme : les spreads du snapshot synthétique
décroissent avec la maturité, ce qui rend la base du forward négative au-delà d'une quinzaine d'années,
un artefact des données et non un comportement de marché.
\begin{figure}[H]
\centering
\begin{tikzpicture}
\begin{groupplot}[group style={group size=2 by 1, horizontal sep=1.8cm}, width=0.48\linewidth,
                  height=6cm]
\nextgroupplot[xlabel={Maturité (années)}, ylabel={Taux zéro (\%)}, legend pos=south east]
\addplot[blue!70!black, thick] table[x=t, y=ois]{data/irs_curves.dat};
\addplot[orange!85!black, thick] table[x=t, y=e6]{data/irs_curves.dat};
\addplot[green!50!black, thick, dashed] table[x=t, y=e3]{data/irs_curves.dat};
\legend{OIS, Euribor 6M, Euribor 3M}
\nextgroupplot[xlabel={Date de début du forward 6M (années)}, ylabel={Forward 6M Euribor $-$ OIS (pb)}]
\addplot[red!70!black, thick] table[x=t, y=basis]{data/irs_curves.dat};
\end{groupplot}
\end{tikzpicture}
\caption{Courbes de taux zéro continus (Act/365F) et base entre le forward 6M de la courbe Euribor 6M
et celui de la courbe OIS.}\label{fig:curves}
\end{figure}

\subsection{Calibration de la courbe Euribor 6M}
\IrsPillarCount{} instruments (un dépôt, quinze IRS), au plus \IrsMaxIterations{} itérations de
Brent par pilier, résidu maximal \IrsMaxResidual{} sur le taux.
\begin{table}[H]
\centering
\small
\begin{tabular}{@{}llcrrrcr@{}}
\toprule
Pilier & Type & Date & Cotation (\%) & DF & Zéro (\%) & Itér. & Résidu \\
\midrule
6M & Dépôt & 2027-04-21 & 2{,}014 & 0{,}98981054 & 2{,}032 & 9 & $9{,}4\times10^{-17}$ \\
1Y & IRS & 2027-10-21 & 2{,}035 & 0{,}97994165 & 2{,}015 & 9 & $-4{,}5\times10^{-17}$ \\
2Y & IRS & 2028-10-23 & 2{,}157 & 0{,}95796499 & 2{,}133 & 9 & $-1{,}7\times10^{-16}$ \\
3Y & IRS & 2029-10-22 & 2{,}270 & 0{,}93460852 & 2{,}246 & 9 & $-5{,}9\times10^{-17}$ \\
4Y & IRS & 2030-10-21 & 2{,}359 & 0{,}91062619 & 2{,}336 & 9 & $-6{,}2\times10^{-17}$ \\
5Y & IRS & 2031-10-21 & 2{,}432 & 0{,}88631064 & 2{,}410 & 9 & $-6{,}9\times10^{-17}$ \\
6Y & IRS & 2032-10-21 & 2{,}495 & 0{,}86187006 & 2{,}473 & 9 & $-5{,}9\times10^{-17}$ \\
7Y & IRS & 2033-10-21 & 2{,}553 & 0{,}83729396 & 2{,}533 & 9 & $6{,}9\times10^{-17}$ \\
8Y & IRS & 2034-10-23 & 2{,}606 & 0{,}81263478 & 2{,}588 & 9 & $-5{,}2\times10^{-17}$ \\
9Y & IRS & 2035-10-22 & 2{,}653 & 0{,}78841209 & 2{,}637 & 9 & $-4{,}2\times10^{-17}$ \\
10Y & IRS & 2036-10-21 & 2{,}695 & 0{,}76453129 & 2{,}681 & 9 & $-3{,}5\times10^{-17}$ \\
12Y & IRS & 2038-10-21 & 2{,}760 & 0{,}71859002 & 2{,}751 & 10 & $-1{,}0\times10^{-17}$ \\
15Y & IRS & 2041-10-21 & 2{,}818 & 0{,}65554538 & 2{,}812 & 10 & $3{,}5\times10^{-18}$ \\
20Y & IRS & 2046-10-22 & 2{,}836 & 0{,}56795072 & 2{,}826 & 13 & $-3{,}5\times10^{-18}$ \\
25Y & IRS & 2051-10-23 & 2{,}787 & 0{,}50186338 & 2{,}755 & 12 & $-1{,}7\times10^{-17}$ \\
30Y & IRS & 2056-10-23 & 2{,}721 & 0{,}44975590 & 2{,}661 & 11 & $0$ \\
\bottomrule
\end{tabular}

\caption{Nœuds calibrés de la courbe Euribor 6M.}
\end{table}

\subsection{Flux et valeur actuelle}
Le swap a \IrsFixedCount{} coupons fixes et \IrsFloatingCount{} coupons flottants. Le \emph{premier}
coupon flottant, fixé le \IrsFirstFixingDate{} (la date de valorisation), prend le fixing publié de
\pct{\IrsFirstFixing} au lieu d'un forward projeté ; les autres sont projetés sur la courbe Euribor 6M
et tous sont actualisés sur la courbe OIS.
\begin{table}[H]
\centering
\scriptsize
\resizebox{\linewidth}{!}{\begin{tabular}{@{}lcccrrrr@{}}
\toprule
Jambe & Début & Fin & Paiement & Taux (\%) & Montant & DF OIS & PV \\
\midrule
Fixe & 2026-10-21 & 2027-10-21 & 2027-10-21 & 2{,}5000 & $-$250\,000{,}00 & 0{,}981142 & $-$245\,285{,}47 \\
Fixe & 2027-10-21 & 2028-10-23 & 2028-10-23 & 2{,}5000 & $-$251\,388{,}89 & 0{,}960519 & $-$241\,463{,}74 \\
Fixe & 2028-10-23 & 2029-10-22 & 2029-10-22 & 2{,}5000 & $-$249\,305{,}56 & 0{,}938515 & $-$233\,976{,}92 \\
Fixe & 2029-10-22 & 2030-10-21 & 2030-10-21 & 2{,}5000 & $-$249\,305{,}56 & 0{,}915814 & $-$228\,317{,}41 \\
Fixe & 2030-10-21 & 2031-10-21 & 2031-10-21 & 2{,}5000 & $-$250\,000{,}00 & 0{,}892679 & $-$223\,169{,}63 \\
Flottante & 2026-10-21 & 2027-04-21 & 2027-04-21 & 2{,}0140 & 101\,818{,}89 & 0{,}990509 & 100\,852{,}53 \\
Flottante & 2027-04-21 & 2027-10-21 & 2027-10-21 & 1{,}9812 & 100\,709{,}02 & 0{,}981142 & 98\,809{,}84 \\
Flottante & 2027-10-21 & 2028-04-21 & 2028-04-21 & 2{,}2314 & 113\,431{,}13 & 0{,}971114 & 110\,154{,}58 \\
Flottante & 2028-04-21 & 2028-10-23 & 2028-10-23 & 2{,}2316 & 114\,677{,}89 & 0{,}960519 & 110\,150{,}27 \\
Flottante & 2028-10-23 & 2029-04-23 & 2029-04-23 & 2{,}4563 & 124\,182{,}14 & 0{,}949455 & 117\,905{,}35 \\
Flottante & 2029-04-23 & 2029-10-22 & 2029-10-22 & 2{,}4563 & 124\,182{,}14 & 0{,}938515 & 116\,546{,}76 \\
Flottante & 2029-10-22 & 2030-04-23 & 2030-04-23 & 2{,}5878 & 131\,548{,}24 & 0{,}927034 & 121\,949{,}68 \\
Flottante & 2030-04-23 & 2030-10-21 & 2030-10-21 & 2{,}5876 & 130\,101{,}24 & 0{,}915814 & 119\,148{,}48 \\
Flottante & 2030-10-21 & 2031-04-21 & 2031-04-21 & 2{,}6875 & 135\,868{,}87 & 0{,}904205 & 122\,853{,}32 \\
Flottante & 2031-04-21 & 2031-10-21 & 2031-10-21 & 2{,}6876 & 136\,620{,}48 & 0{,}892679 & 121\,958{,}17 \\
\bottomrule
\end{tabular}
}
\caption{Flux du swap (montants signés : positif $=$ reçu, négatif $=$ payé).}
\end{table}
\begin{center}
\begin{tabular}{@{}lr@{}}
\toprule
Grandeur & Valeur \\
\midrule
PV de la jambe fixe (payée) & \IrsFixedLegPV\eur \\
PV de la jambe flottante (reçue) & \IrsFloatingLegPV\eur \\
\textbf{PV du swap payeur} & \textbf{\IrsPV\eur} \\
Annuité \(A\) & \IrsAnnuity\eur \\
PVBP \(=A\times10^{-4}\) & \IrsPVBP\eur \\
Taux de par \(S\) & \pct{\IrsParRate} \\
Écart au par \(S-K\) & \IrsSpread{} pb \\
\bottomrule
\end{tabular}
\end{center}
On retrouve l'identité \eqref{eq:par} : \(\PV=A\,(S-K)\). L'annuité multipliée par l'écart au par
de \IrsSpread{} pb (soit \IrsSpread{} fois \(10^{-4}\) en valeur décimale) redonne la valeur actuelle
du tableau.

\subsection{Ce que change le multi-courbes}
Si l'on projetait par erreur l'Euribor 6M sur la courbe OIS (cadre mono-courbe), le même swap vaudrait
\IrsSingleCurvePV\eur, soit un écart de \IrsSingleCurveGap\eur{} : sur ce marché synthétique, la
prime de l'Euribor 6M sur l'OIS pèse bien plus que le PVBP. C'est l'ordre de grandeur de l'erreur que
le cadre multi-courbes supprime.

\subsection{Risque}
\begin{center}
\begin{tabular}{@{}lr@{}}
\toprule
Mesure & Valeur \\
\midrule
DV01 parallèle & \IrsParallelDV\eur/pb \\
DV01 courbe Euribor 6M (projection) & \IrsDVSix\eur/pb \\
DV01 courbe OIS (actualisation) & \IrsDVOis\eur/pb \\
DV01 courbe Euribor 3M & \IrsDVThree\eur/pb \\
Somme des DV01 par instrument & \IrsKeyRateSum\eur/pb (écart au parallèle : \IrsKeyRateGap{}) \\
Convexité & \IrsConvexity{} \eur/pb\(^2\) \\
Nombre de revalorisations & \IrsRiskBumps{} \\
\bottomrule
\end{tabular}
\end{center}
Le risque se concentre sur la courbe de projection Euribor 6M : un swap payeur à 5 ans est surtout
exposé à l'indice qu'il reçoit ; l'actualisation OIS pèse très peu et le swap n'a aucune exposition à
l'Euribor 3M. La figure~\ref{fig:keyrate} détaille le DV01 par instrument : l'essentiel est sur le
swap Euribor 6M à 5 ans (son DV01 vaut le PVBP, puisque la cotation 5Y est le taux de par), et
l'exposition s'annule au-delà de l'horizon du swap (maturités longues), exactement et non par
approximation. Le bucket du dépôt Euribor 6M est \emph{négatif} : relever ce dépôt, toutes les autres
cotations restant fixes, abaisse \(P^e\) à six mois et donc le forward 6M--1Y, qui entre dans la jambe
reçue. Ce bucket compense en partie le bucket 5Y, d'où un DV01 parallèle inférieur au PVBP.
\begin{figure}[H]
\centering
\begin{tikzpicture}
\begin{axis}[ybar, bar width=4pt, width=\linewidth, height=6.2cm, xtick=data,
  xticklabels from table={data/irs_keyrate.dat}{label},
  x tick label style={rotate=60, anchor=east, font=\scriptsize},
  ylabel={DV01 (\eur/pb)}, enlarge x limits=0.03, legend pos=north west]
\addplot[fill=orange!70] table[x=i, y=e6]{data/irs_keyrate.dat};
\addplot[fill=blue!60] table[x=i, y=ois]{data/irs_keyrate.dat};
\legend{Euribor 6M, OIS}
\end{axis}
\end{tikzpicture}
\caption{DV01 par instrument de calibration du swap de l'exemple.}\label{fig:keyrate}
\end{figure}

\paragraph{Contrôle de Taylor.} Un choc parallèle de \IrsTaylorShift{} pb donne, après re-calibration,
une variation de PV de \IrsTaylorActual\eur, contre \IrsTaylorEstimate\eur{} prédits par
\(\mathrm{DV01}\,x+\tfrac12\mathrm{convexité}\,x^2\) : le DV01 et la convexité sont cohérents.

\section{Validation contre QuantLib}\label{sec:valid}

QuantLib est l'oracle de test (jamais importé par le cœur). La comparaison se fait couche par couche,
sur des grilles générées (9 marchés : 3 dates de valorisation, dont un mois-fin et une veille de
Pâques, et 3 chocs parallèles de \(-250\), \(0\) et \(+150\) pb sur toutes les cotations, Euribor
comprises). Les courbes QuantLib sont figées avant valorisation pour que les fixings chargés dans
l'indice ne les re-calibrent pas (\S\ref{sec:nofixing}).
\begin{center}
\begin{tabularx}{\linewidth}{@{}Xrr@{}}
\toprule
Comparaison & Écart maximal & Critère \\
\midrule
Facteurs d'actualisation, courbes Euribor 3M et 6M (log-linéaire) & \(3{,}5\times10^{-13}\) & \(10^{-10}\) \\
Idem, interpolation linéaire sur les zéros & \(3{,}0\times10^{-12}\) & \(10^{-10}\) \\
PV d'IRS par notionnel (972 swaps spot + 72 en cours / forward) & \(2{,}2\times10^{-13}\) & \(10^{-8}\) \\
Taux de par, relatif & \(4{,}0\times10^{-10}\) & \(10^{-8}\) \\
Annuité, relative & \(5{,}6\times10^{-14}\) & \(10^{-8}\) \\
Échéanciers (accrual, paiement, fixing, 2 jambes, stubs 18M/30M) & égalité exacte & exacte \\
DV01 par instrument, rapporté au DV01 parallèle (3 trades, 56 buckets) & \(9{,}8\times10^{-12}\) & \(10^{-6}\) \\
DV01 parallèle, relatif & \(8{,}1\times10^{-13}\) & \(10^{-6}\) \\
Somme des buckets contre DV01 parallèle, relatif & \(2{,}1\times10^{-7}\) & \(10^{-6}\) \\
\bottomrule
\end{tabularx}
\end{center}
Le DV01 par instrument est comparé au DV01 \emph{parallèle du trade} : un instrument loin de l'échéance a
un DV01 proche de zéro, pour lequel un écart relatif à lui-même n'aurait pas de sens. Les comparaisons de
risque utilisent le pricer de production (courbes tronquées, mémoïsation) face aux courbes complètes de
QuantLib, ce qui valide aussi les optimisations.

Les autres garde-fous : propriété \og un swap au taux de par vaut zéro \fg, symétrie payeur/receveur,
propriété de somme des DV01 (\code{hypothesis}), égalité au bit près entre une reconstruction mémoïsée et
une reconstruction à neuf, test d'intégration SDK $\to$ HTTP $\to$ cœur et \emph{golden file} de la
phase~2.

\section{Limites et prolongements}
\begin{itemize}
  \item Données synthétiques : les niveaux de taux et de base ne sont pas ceux du marché ; le fixing
    du jour de l'Euribor est pris égal au dépôt spot du snapshot.
  \item Pas de FRA, de \emph{basis swap} ni de swap avec stub à la cotation (la maturité 18M n'est pas
    calibrée sur l'Euribor).
  \item Risque par bump-and-reprice (différences finies) ; la différentiation automatique, hors v1,
    réduirait le coût sans changer le contrat.
  \item Phase~3 : scénarios (parallèle, \emph{steepener}, \emph{flattener}, \emph{twist}), carry et
    \emph{roll-down}, USD (SOFR) par simple configuration.
\end{itemize}

\appendix
\section{Reproduire le document}
\begin{lstlisting}[language=bash,caption={Régénérer les chiffres puis compiler.},numbers=none]
uv run python docs/latex/generate_irs_data.py     # nombres, tableaux, figures
pdflatex pricing_irs.tex && pdflatex pricing_irs.tex   # deux passes
\end{lstlisting}
Le script appelle le pricer sur le snapshot \code{EUR_2026-10-19} et écrit \texttt{generated/irs\_*.tex}
et \texttt{data/irs\_*.dat} : aucun nombre de ce document n'est saisi à la main, à l'exception des écarts
de validation (tableau de la \S\ref{sec:valid}), mesurés par les tests de \code{tests/validation/} et
consignés dans l'ADR-031.

\end{document}
